\chapter{Introduction}
\label{cha-region-intro}

This document describes the \FramaC/\Region plug-in that creates a map
dispatching each memory location into regions, modularly for each \C function.
The \Region plug-in aims at producing a correct, yet optimistic, map of the
memory. By optimistic, we mean that the plug-in assumes that the programmer
tries to maximize separation between pointers. Of course, that might not be
true, thus the plug-in can show what has been computed so that one can provide
additional annotations to tell the plug-in that some pointer might be in alias.

While the plug-in can be run manually, using the \lstinline{-region} option of
\FramaC. Its main purpose is to support other analyzers through its API.

The document is structured as follows:
\begin{itemize}
  \item Section~\ref{cha-region-plugin} describes the meaning of memory regions
        and how these regions are computed from input source code,
  %\item Section~\ref{region_logic} ...
  %\item Section~\ref{region_api} ...
\end{itemize}
